بردار تصادفی $x \sim \mathcal{N}(0,\Sigma)$ با
$\Sigma=\begin{bmatrix}1&0\\0&4\end{bmatrix}$
و $x=\begin{bmatrix}1\\2\end{bmatrix}$ داده شده است. ابتدا کمیت‌های لازم را حساب می‌کنیم:
\begin{flalign*}
k&=2, \qquad \det(\Sigma)=1\cdot 4=4, \qquad \Sigma^{-1}=\begin{bmatrix}1&0\\0&\frac{1}{4}\end{bmatrix} &&
\end{flalign*}
\begin{flalign*}
(x-\mu)^{T}\Sigma^{-1}(x-\mu)
&=\begin{bmatrix}1&2\end{bmatrix}
\begin{bmatrix}1&0\\0&\frac{1}{4}\end{bmatrix}
\begin{bmatrix}1\\2\end{bmatrix}
=1+\frac{4}{4}=2 &&
\end{flalign*}
حال تابع چگالی را می‌نویسیم:
\begin{flalign*}
p(x)
&=\frac{1}{\sqrt{(2\pi)^k\det(\Sigma)}}\exp\!\left(-\frac{1}{2}(x-\mu)^T\Sigma^{-1}(x-\mu)\right) \\
&=\frac{1}{\sqrt{(2\pi)^2\cdot 4}}\exp(-1)
=\frac{1}{4\pi}e^{-1} &&
\end{flalign*}
بنابراین مقدار چگالی در $x=[1,2]^T$ برابر است با
$p(x)=\frac{1}{4\pi}e^{-1}$.